\section{Correspondence with the formal statement}\label{sec:correspondence}

This section states the existence theorem in the exact form of the formal statement and deduces it from Theorem~\ref{thm:B}. The formal statement constrains abstract hierarchy data by the conditions Q0--Q6, of which Q1 is Definition~\ref{def:Q1}. They do not impose every property established above: the correlations need not be multilinear, the logarithmic slope $b_*$ is only required to be positive, and the degree-four field~$\Theta$ is not identified as a stress trace. Section~\ref{sec:formal-conventions} records the conventions in which the formal statement differs from the main text.

\subsection{Hierarchy data and the conditions Q0--Q6}\label{sec:hierarchy-data}

\begin{definition}[Q0. Hierarchy data]\label{def:Q0}
Hierarchy data over a Wilson family~$\mathcal W$ consist of:
\begin{itemize}[nosep]
\item complex Hilbert and vector spaces~$\mathcal H,\mathcal L$, and a relativistic datum with the group of Remark~\ref{rem:formal-poincare};
\item finite-dimensional complex normed spaces~$F_d$, linear maps~$\iota_d:F_d\to\mathcal L$, and $\pi_{\rm quo}:\mathcal T\to\mathbb F$, where $\mathbb F=\bigsqcup_dF_d$;
\item a compact metric space~$\phases$, a Borel measure~$\nu$ and $\xi_*\in\phases$;
\item functions $S_n^\zeta:\mathbb F^n\times(\R^4)^n\to\C$ for $\zeta\in\phases\sqcup\{\mathrm{mix}\}$, and $\mathsf F_n:\mathbb F^n\times\C^n\times(\R^3)^n\to\C$;
\item finite sets~$\channels{j}$, fields~$O_{j,\lambda}\in F_j$, and coefficients $c_{j,\lambda}:\mathbb F^{I_{\rm cl}}\times(\R^4)^{I_{\rm cl}}\times\R^4\to\C$ for every finite index set~$I_{\rm cl}$;
\item an element~$\Theta\in F_4$.
\end{itemize}
\end{definition}

The spaces, maps and functions of Q0 are arbitrary on the displayed domains; put $\deg(d,A)=d$ and $\iota(d,A)=\iota_dA$. The conditions below constrain the correlations at separated positions. In Q2--Q6 all cited definitions and formulas are read with these hierarchy data and with the conventions of Section~\ref{sec:formal-conventions}; in definitions indexed by a phase, read $\Omega_{\xi_*}$, $\Phi_{\xi_*}$ and $\mathsf F_n^{\xi_*}$ as the $\Omega$, $\Phi$ and $\mathsf F_n$ of these data.

\begin{definition}[Q2. Wightman axioms]\label{def:Q2}
The relativistic datum of Q0 is a Wightman theory in the sense of Definition~\ref{def:wightman}.
\end{definition}

\begin{definition}[Q3. Labels and phases]\label{def:Q3}
\textup{(a)} Every nonzero~$\ell\in\mathcal L$ has $\Phi(\ell,f)\ne0$ for some~$f\in\mathcal S$.
\textup{(b)} Each~$\iota_d$ is injective, their images cover~$\mathcal L$, and $\pi_{\rm quo}$ is onto.
\textup{(c)} $\Phi(\iota(A),f)\ne0$ for some~$A\in\mathbb F$ and $f\in\mathcal S$.
\textup{(d)} The measure~$\nu$ is a probability measure of full support, and for every $n$, $\boldsymbol A\in\mathbb F^n$ and $x\in\Conf{n}$, $\xi\mapsto S_n^\xi(\boldsymbol A;x)$ is continuous and $S_n^{\rm mix}(\boldsymbol A;x)=\int_{\phases}S_n^\xi(\boldsymbol A;x)\,d\nu(\xi)$.
\textup{(e)} Each~$S_n^\zeta$ is simultaneously permutation symmetric and smooth on~$\Conf{n}$.
\end{definition}

\begin{definition}[Q4. Selected continuation]\label{def:Q4}
The functions~$\mathsf F_n$ satisfy the ordered-time continuation condition of Definition~\ref{def:D15} and the boundary condition of Definition~\ref{def:D14f} at~$\xi_*$, read as in Remark~\ref{rem:integral-zero}, using the relativistic datum and the maps~$\iota_d$ of the hierarchy data.
\end{definition}

\begin{definition}[Q5. Products]\label{def:Q5}
For the channels, fields, coefficients and tables~$S^\zeta$ of Q0, form~$\mathcal R_D^\zeta$ by Definition~\ref{def:D18}. These coefficients and remainders satisfy the two inequalities of Theorem~\ref{prop:OP}: the coefficient threshold depends on $I_{\rm cl},\boldsymbol A,M_{\rm geo},\eta,\varepsilon,j,\lambda$; the remainder threshold depends on $I_{\rm cl},I_{\rm sp},\boldsymbol A,\boldsymbol B,M_{\rm geo},\eta,\delta_{\rm sp},\varepsilon,D$. Each threshold is uniform in positions and state indices.
\end{definition}

\begin{definition}[Q6. Trace test]\label{def:Q6}
For the chosen~$\Theta\in F_4$, with $\tracetwopoint$ its time-axis two-point function in the phase~$\xi_*$ (Section~\ref{sec:short-distance-defs}) and $\traceamplitude=3\dimG/\pi^4$, require $\langle\Omega,\Phi(\iota_4\Theta,f)\Omega\rangle=0$ for every $f\in\mathcal S$,
\begin{equation*}
\lim_{r\downarrow0}r^8\log^2(1/r)\tracetwopoint(r)=\traceamplitude>0
\quad\text{in }\C\,,
\end{equation*}
and the finite-profile exclusion of Proposition~\ref{prop:NF}.
\end{definition}

For the lattice limit, $S^{\rm mix}\circ\pi_{\rm quo}$ denotes the kernels $K_n(v;x)=S_n^{\rm mix}(\pi_{\rm quo}(v_1),\ldots,\pi_{\rm quo}(v_n);x)$.

\setcounter{maintheorem}{0}
\begin{maintheorem}[Existence]\label{thm:A}
For every gauge group of Definition~\ref{def:D01}, there is a regular \textup{(Q1)} Wilson family~$\mathcal W$ such that, for every strictly increasing~$\mathsf p:\N\to\N$, there are a strictly increasing~$\mathsf q$ and hierarchy data over~$\mathcal W$ \textup{(Q0)} satisfying \textup{Q2--Q6}, with
\begin{equation*}
\mathcal W\longrightarrow S^{\rm mix}\circ\pi_{\rm quo}
\quad\text{along }P=\mathsf p\circ\mathsf q
\quad\text{in the sense of Definition~\ref{def:D09}}\,.
\end{equation*}
\end{maintheorem}

Theorem~\ref{thm:A} fixes the lattice theory before the volume sequence. Q2 gives relativistic fields; Q3 and Q4 express the limiting correlations as a phase average and continue the selected phase to those fields; Q5 approximates local products uniformly in the state and in separated spectators; Q6 fixes the vacuum centering and the short-distance law of the distinguished degree-four field.

\begin{proof}[Deduction of Theorem~\ref{thm:A}]
\begin{enumerate}[label=\textup{\arabic*.},ref=\arabic*,start=10]
\item\label{step:deduction-conclusion} Steps~\ref{step:deduction-choice}--\ref{step:deduction-products} of Section~\ref{sec:deduction} give Theorem~\ref{thm:B} and hierarchy data satisfying Q2--Q6 on the same family and volume refinement. Fix any~$\blockpenalty>0$, for instance~$\blockpenalty=1$. The regular family precedes~$\mathsf p$; its refinement and hierarchy data follow it, with the ordered convergence of Step~\ref{step:deduction-hierarchy}. This proves Theorem~\ref{thm:A}.\qedhere
\end{enumerate}
\end{proof}

As in the formal statement, Q0--Q3 do not require $S_n^\zeta$ to be multilinear in the fields; in the construction they are, since $S$ and every $S^\xi$ are tables (Theorems~\ref{prop:IV} and~\ref{prop:PH}). Theorem~\ref{thm:A} leaves $b_*>0$ unspecified; Theorem~\ref{prop:MT} supplies $b_*=\slope=2b_0=11\casimir/(24\pi^2)$.

\begin{remark}[Nontriviality]\label{rem:nontriviality}
Q3(c) requires a nonzero smeared field; given Q3(a)--(b) this is equivalent to $\mathcal L\ne0$, and in the construction it is met by the unit field (Step~\ref{step:deduction-hierarchy}). Nontriviality comes from the trace: $\Theta\ne0$ (Theorem~\ref{prop:TR}), so faithfulness gives $\Phi(\iota_4\Theta,f)\ne0$ for some~$f$, while its vacuum expectation vanishes; hence it is not a multiple of the identity. Q6 gives the positive limit $3\dimG/\pi^4$ of $r^8\log^2(1/r)\tracetwopoint(r)$, and on no interval $(0,\delta)$ does $\tracetwopoint$ agree with a finite Wick profile. Thus every Gaussian or finite free-field realization whose two-point profile for~$\Theta$ belongs to the class of Definition~\ref{def:D17} is excluded.
\end{remark}

\subsection{Formal conventions}\label{sec:formal-conventions}

The formal statement differs from the main text in the following conventions; none of them changes the content of the results.

\begin{remark}[Abstract Poincar\'e group]\label{rem:formal-poincare}
In the formal statement a relativistic datum carries a set $\mathcal P'$ with a unit, a product and a topology, and a homeomorphism $\kappa_{\rm P}:\mathcal P'\to SO^+(1,3)\ltimes\R^4$ respecting the identity and the products (which makes $\mathcal P'$ a topological group), a strongly continuous unitary representation $\mathsf U$ of $\mathcal P'$ with $\mathsf U(\kappa_{\rm P}^{-1}(1,z))=\mathsf T(z)$, label maps $\rho_{\rm lab}$ on~$\mathcal P'$, and the action on tests through $\kappa_{\rm P}$. Transport along $\kappa_{\rm P}$ turns such a datum into one for the ordinary group of Definition~\ref{def:D14a}, and conversely, and it preserves W0--W3. Definition~\ref{def:wightman} is therefore equivalent to the formal one.
\end{remark}

\begin{remark}[Nonintegrable integrands]\label{rem:integral-zero}
In the boundary condition of Q4 the formal statement does not require integrability: there a nonintegrable function has integral zero. For the continuations constructed in Theorem~\ref{prop:RC} the integrands are integrable for every $r>0$ (by the bound~\eqref{eq:H-poly} and rapid decay of the tests), so this convention is never used for them, and Definition~\ref{def:D14f} requires integrability.
\end{remark}

The formal development also states the volume limit as local uniform convergence, on $\Gen{n}$, of the fixed-volume kernels $M^{\rm UV}_{P(k),v}$ restricted to the periodic generic domains $\{x:x_i^a-x_j^a\notin L_{P(k)}\Z\ (i\ne j,\ a=1,\ldots,4)\}$. This form follows from~\eqref{eq:F-conf}: a compact subset of~$\Gen{n}$ has bounded diameter and positive minimal coordinate separation, so it eventually lies in the periodic generic domain, and $\Gen{n}\subset\Conf{n}$.

\subsection{An averaged countertest}\label{sec:averaged-countertest}

The averaged Fock countertest of Section~\ref{sec:partII-I} belongs to this correspondence: it compares the trace field with free comparison vectors modified by arbitrary maps that preserve physical inner products and fix $U\mathcal H$. Its integrands are constant, so the convention of Remark~\ref{rem:integral-zero} is not needed for them.

Let $n\ge0$ and let $U$ be a comparison (Definition~\ref{def:D30}), with $\mathcal F_n$ as in Definition~\ref{def:D29}. For a probability space $(Z,\lambda)$, let each $P_z:\mathcal F_n\to\mathcal F_n$ preserve physical inner products and fix $U\mathcal H$ pointwise, without assuming linearity.
The maps $P_z$ model transformations of the free comparison vectors that fix every embedded interacting state in $U\mathcal H$. Because physical inner products are preserved, neither matrix element in Lemma~\ref{lem:I2.3} changes, so the averaged inequality states that averaging these transformations cannot restore the relative-locality identity.
\begin{lemma}[Localized averaged countertest]\label{lem:I2.4}
For every such $n,U,(Z,\lambda),(P_z)$ and $r>0$, there are open $A,B$, with $\overline A\subset B(0,r)$ and $\minksquare(x-y)<0$ on $\overline A\times\overline B$, and tests $f,h$ as in Lemma~\ref{lem:I2.3}, supported respectively in $A,B$, such that
\begin{equation*}
\int_Z\langle P_z\mathcal E(-h),\psi_U(f)\rangle\,d\lambda(z)
\ne\int_Z\langle\psi_U(\bar f),P_z\mathcal E(h)\rangle\,d\lambda(z)\,.
\end{equation*}
\end{lemma}
\begin{proof}
For physical $v$ and $w\in U\mathcal H$ (Definition~\ref{def:D30}), $\langle P_zv,w\rangle=\langle P_zv,P_zw\rangle=\langle v,w\rangle$.
Thus averaging preserves both integrands.

Apply Lemma~\ref{lem:I2.3} in $B(0,r/2)$. Compactness and strict spacelike separation give support neighbourhoods $A,B$ with spacelike separated closures and $\overline A\subset B(0,r)$.
\end{proof}
